\documentclass[12pt, a4paper]{article}

% ── Packages ──────────────────────────────────────────────────────────────────
\usepackage[T1]{fontenc}
\usepackage[utf8]{inputenc}
\usepackage[french]{babel}
\usepackage[margin=2.4cm, top=2.8cm, bottom=2.8cm]{geometry}
\usepackage{lmodern}
\usepackage{microtype}
\usepackage{parskip}
\usepackage{titlesec}
\usepackage{xcolor}
\usepackage{mdframed}
\usepackage{enumitem}
\usepackage{booktabs}
\usepackage{array}
\usepackage{fancyhdr}
\usepackage{hyperref}
\usepackage{tikz}
\usepackage{pgfplots}
\pgfplotsset{compat=1.18}
\usetikzlibrary{shapes.geometric, shapes.rounded, arrows.meta, positioning,
                fit, backgrounds, shadows.blur, calc}

% ── Colors ────────────────────────────────────────────────────────────────────
\definecolor{accent}{RGB}{94, 106, 210}
\definecolor{accentlight}{RGB}{220, 223, 245}
\definecolor{lightgray}{RGB}{245, 246, 248}
\definecolor{darkgray}{RGB}{50, 55, 70}
\definecolor{mutedtext}{RGB}{110, 115, 135}
\definecolor{csource}{RGB}{167, 139, 250}
\definecolor{cengine}{RGB}{94, 106, 210}
\definecolor{coutput}{RGB}{52, 211, 153}
\definecolor{cdanger}{RGB}{239, 68, 68}
\definecolor{cokay}{RGB}{34, 197, 94}
\definecolor{camber}{RGB}{245, 158, 11}

% ── Hyperref ──────────────────────────────────────────────────────────────────
\hypersetup{colorlinks=true, linkcolor=accent, urlcolor=accent,
            pdftitle={dedomena -- Presentation Produit}}

% ── Section styles ────────────────────────────────────────────────────────────
\titleformat{\section}
  {\large\bfseries\color{darkgray}}{}{}{}
  [\vspace{-0.3em}\textcolor{accent}{\rule{\linewidth}{1.5pt}}\vspace{0.4em}]
\titleformat{\subsection}{\normalsize\bfseries\color{darkgray}}{}{}{}
\titlespacing{\section}{0pt}{1.8em}{0.6em}
\titlespacing{\subsection}{0pt}{1.2em}{0.3em}

% ── Header / footer ───────────────────────────────────────────────────────────
\pagestyle{fancy}
\fancyhf{}
\fancyhead[L]{\small\textcolor{mutedtext}{dedomena \textsigma}}
\fancyhead[R]{\small\textcolor{mutedtext}{Presentation produit}}
\fancyfoot[C]{\small\textcolor{mutedtext}{\thepage}}
\renewcommand{\headrulewidth}{0pt}

% ── Callout boxes ─────────────────────────────────────────────────────────────
\newmdenv[backgroundcolor=lightgray, linecolor=accent, linewidth=3pt,
          topline=false, bottomline=false, rightline=false,
          innerleftmargin=14pt, innerrightmargin=14pt,
          innertopmargin=10pt, innerbottommargin=10pt,
          skipabove=10pt, skipbelow=10pt]{callout}

\newmdenv[backgroundcolor=accent!7, linecolor=accent!35, linewidth=1pt,
          roundcorner=5pt, innerleftmargin=14pt, innerrightmargin=14pt,
          innertopmargin=10pt, innerbottommargin=10pt,
          skipabove=10pt, skipbelow=10pt]{highlight}

% ── TikZ styles ───────────────────────────────────────────────────────────────
\tikzset{
  box/.style={
    rectangle, rounded corners=6pt,
    draw=#1!60, fill=#1!10, text=#1!80!black,
    minimum height=1.0cm, text centered,
    font=\small\bfseries, inner sep=8pt,
    drop shadow={shadow xshift=1pt, shadow yshift=-1pt,
                 fill=#1!20, opacity=0.5}
  },
  arr/.style={-{Stealth[length=6pt]}, thick, color=#1,
              shorten >=2pt, shorten <=2pt},
  alabel/.style={font=\scriptsize\itshape, text=mutedtext,
                 fill=white, inner sep=2pt, midway, sloped=false},
}

% ── Document ──────────────────────────────────────────────────────────────────
\begin{document}

% ── Page de couverture ────────────────────────────────────────────────────────
\thispagestyle{empty}
\begin{center}
  \vspace*{1.4cm}
  {\fontsize{36}{42}\selectfont\bfseries\color{darkgray} dedomena}
  \hspace{0.4em}
  {\fontsize{36}{42}\selectfont\color{accent!45}\textsigma}

  \vspace{0.6em}
  {\large\color{mutedtext} Une facon plus intelligente de travailler avec vos donnees}

  \vspace{1.6em}
  \textcolor{accent!30}{\rule{7cm}{0.6pt}}

  \vspace{1em}
  {\small\color{mutedtext}
    Presentation produit \quad\textbullet\quad Document interne \quad\textbullet\quad \today
  }
  \vspace{0.6cm}

  % Diagramme de couverture
  \begin{tikzpicture}[node distance=0.8cm and 1.1cm]
    \node[box=csource, minimum width=3.2cm] (src) {Vos donnees};
    \node[box=cengine, minimum width=3.2cm, right=of src] (eng) {dedomena};
    \node[box=coutput, minimum width=3.2cm, right=of eng] (out) {Reponse claire};

    \draw[arr=csource] (src) -- (eng)
      node[alabel, above] {connecte};
    \draw[arr=cengine] (eng) -- (out)
      node[alabel, above] {livre};

    \node[font=\scriptsize\color{mutedtext}, below=0.35cm of src]
      {100+ sources};
    \node[font=\scriptsize\color{mutedtext}, below=0.35cm of eng]
      {raisonnement IA};
    \node[font=\scriptsize\color{mutedtext}, below=0.35cm of out]
      {en quelques secondes};
  \end{tikzpicture}

  \vspace{1.6cm}
\end{center}
\newpage

% ── 1. Qu'est-ce que dedomena ─────────────────────────────────────────────────
\section{Qu'est-ce que dedomena?}

dedomena est une plateforme web qui permet a n'importe qui dans l'organisation
de poser des questions sur les donnees de l'entreprise en francais courant et
d'obtenir des reponses claires et sourcees en quelques secondes.

Pas de SQL. Pas de formules Excel. Pas d'attente pour qu'un analyste produise
un rapport. Vous connectez vos sources de donnees une seule fois, vous tapez
votre question, et la plateforme fait le reste.

\begin{callout}
  \textbf{Pensez-y comme si vous aviez un assistant tres competent qui a lu
  tous les tableaux, toutes les bases de donnees et tous les rapports de
  l'entreprise, et qui est toujours disponible pour repondre instantanement
  a n'importe quelle question.}
\end{callout}


% ── 2. Le probleme ────────────────────────────────────────────────────────────
\section{Le probleme que nous resolvons}

\subsection{L'ancienne methode est trop lente}

Obtenir une reponse a partir des donnees implique aujourd'hui trop d'etapes
et trop de personnes. Le diagramme ci-dessous montre ce a quoi ressemble
une demande typique en ce moment.

\vspace{0.6em}
\begin{center}
\begin{tikzpicture}[node distance=0.5cm and 0.7cm,
    step/.style={rectangle, rounded corners=4pt, draw=cdanger!50,
                 fill=cdanger!8, text=cdanger!80!black, font=\small,
                 minimum height=0.85cm, minimum width=3.2cm,
                 text centered, inner sep=6pt},
    wait/.style={font=\scriptsize\itshape\color{mutedtext},
                 fill=white, inner sep=2pt}]

  \node[step] (s1) {Le directeur a besoin d'un chiffre};
  \node[step, right=of s1] (s2) {Il contacte l'analyste};
  \node[step, right=of s2] (s3) {L'analyste cherche les donnees};
  \node[step, below=of s3]  (s4) {L'analyste redige une requete};
  \node[step, left=of s4]   (s5) {Il formate un rapport};
  \node[step, left=of s5]   (s6) {Le directeur recoit la reponse};

  \draw[arr=cdanger] (s1)--(s2);
  \draw[arr=cdanger] (s2)--(s3) node[wait, above]{1 a 2 jours};
  \draw[arr=cdanger] (s3)--(s4);
  \draw[arr=cdanger] (s4)--(s5) node[wait, above]{plusieurs heures};
  \draw[arr=cdanger] (s5)--(s6);

  \node[font=\scriptsize\bfseries\color{cdanger},
        below=0.7cm of s5, xshift=1.5cm]
    {Temps total : 1 a 3 jours pour une seule question};
\end{tikzpicture}
\end{center}
\vspace{0.4em}

C'est trop lent. Cela cree un goulot d'etranglement sur un petit nombre de
personnes techniques, et les decisions sont prises sans les informations
qui existent pourtant deja dans l'organisation.

\subsection{Les donnees sont dispersees partout}

Au-dela de la lenteur, un second probleme est que les donnees sont eparpillees.
Aucun outil ne donne une vue d'ensemble, ce qui signifie que chaque question
transversale necessite un travail manuel important.

\vspace{0.6em}
\begin{center}
\begin{tikzpicture}[
    silo/.style={rectangle, rounded corners=5pt,
                 draw=mutedtext!50, fill=lightgray,
                 font=\small, minimum width=2.4cm, minimum height=0.9cm,
                 text centered, inner sep=6pt, text=darkgray},
    question/.style={ellipse, draw=cdanger!60, fill=cdanger!8,
                     font=\small\itshape, text=cdanger!80!black,
                     inner sep=8pt}]

  \node[silo] (crm)  at (0,0)     {CRM};
  \node[silo] (fin)  at (2.8,0)   {Finances};
  \node[silo] (ops)  at (5.6,0)   {Operations};
  \node[silo] (rh)   at (8.4,0)   {RH};
  \node[silo] (files)at (4.2,-1.6){Tableurs};

  \node[question] (q) at (4.2, -3.2)
    {"Quelle region a cru le plus vite ce trimestre?"};

  \foreach \n in {crm, fin, ops, rh, files}{
    \draw[thick, dashed, color=mutedtext!50] (\n) -- (q);
  }

  \node[font=\scriptsize\color{mutedtext}, below=0.25cm of q]
    {Aucun outil ne peut repondre seul. Quelqu'un doit extraire les donnees de chacun manuellement.};
\end{tikzpicture}
\end{center}
\vspace{0.4em}


% ── 3. Comment nous allons utiliser dedomena ──────────────────────────────────
\section{Comment nous allons utiliser dedomena}

\subsection{La nouvelle methode : une question, une reponse}

Avec dedomena, la meme demande qui prenait des jours ne prend plus que quelques
secondes. Le processus entier se reduit a trois etapes.

\vspace{0.6em}
\begin{center}
\begin{tikzpicture}[node distance=1.0cm and 1.4cm,
    step/.style={rectangle, rounded corners=6pt,
                 draw=cokay!60, fill=cokay!8, text=cokay!80!black,
                 font=\small\bfseries, minimum height=1.0cm,
                 minimum width=3.8cm, text centered, inner sep=8pt,
                 drop shadow={shadow xshift=1pt, shadow yshift=-1pt,
                              fill=cokay!20, opacity=0.4}}]

  \node[step] (a) {Poser une question};
  \node[step, right=of a] (b) {dedomena lit toutes les sources};
  \node[step, right=of b] (c) {Reponse livree};

  \draw[arr=cokay] (a)--(b)
    node[midway, above, font=\scriptsize\itshape\color{mutedtext}]{instantanement};
  \draw[arr=cokay] (b)--(c)
    node[midway, above, font=\scriptsize\itshape\color{mutedtext}]{en secondes};

  \node[font=\scriptsize\bfseries\color{cokay},
        below=0.7cm of b]{Temps total : moins de 30 secondes};
\end{tikzpicture}
\end{center}
\vspace{0.4em}

\subsection{Etape 1 : connecter nos sources de donnees}

Nous connectons les outils et bases de donnees que nous utilisons deja. Il s'agit
d'une configuration unique qui prend quelques minutes par source. Une fois
connectees, dedomena garde les donnees disponibles pour les questions. Rien n'est
envoye a un tiers. Tout reste dans notre environnement.

\vspace{0.8em}
\begin{center}
\begin{tikzpicture}[
    cat/.style={rectangle, rounded corners=5pt,
                draw=csource!50, fill=csource!8,
                font=\small\bfseries, text=csource!80!black,
                minimum width=3.2cm, minimum height=0.85cm,
                text centered, inner sep=8pt},
    eng/.style={rectangle, rounded corners=8pt,
                draw=cengine!60, fill=cengine!12,
                font=\normalsize\bfseries, text=cengine!80!black,
                minimum width=2.6cm, minimum height=3.8cm,
                text centered, inner sep=10pt,
                drop shadow={shadow xshift=1.5pt, shadow yshift=-1.5pt,
                             fill=cengine!25, opacity=0.5}},
    arr/.style={-{Stealth[length=5pt]}, thick, color=csource!60}]

  \node[cat] (crm)   at (0,  2.4) {CRM};
  \node[cat] (db)    at (0,  1.2) {Bases de donnees};
  \node[cat] (files) at (0,  0.0) {Fichiers};
  \node[cat] (cloud) at (0, -1.2) {Stockage cloud};
  \node[cat] (fin)   at (0, -2.4) {Finances};

  \node[eng] (eng) at (4.6, 0) {dedomena};

  \node[rectangle, rounded corners=6pt,
        draw=coutput!60, fill=coutput!10,
        font=\small\bfseries, text=coutput!80!black,
        minimum width=2.8cm, minimum height=1.2cm,
        text centered, inner sep=10pt,
        drop shadow={shadow xshift=1pt, shadow yshift=-1pt,
                     fill=coutput!25, opacity=0.4}]
       (out) at (8.6, 0) {Votre reponse};

  \foreach \s in {crm, db, files, cloud, fin}{
    \draw[arr] (\s) -- (eng);
  }

  \draw[-{Stealth[length=6pt]}, thick, color=cengine!70] (eng) -- (out)
    node[midway, above, font=\scriptsize\itshape\color{mutedtext}]{en secondes};

  \node[font=\scriptsize\color{mutedtext}, below=0.3cm of fin]
    {Plus de 100 sources supportees};
\end{tikzpicture}
\end{center}
\vspace{0.4em}

\subsection{Etape 2 : poser des questions en langage naturel}

N'importe quel membre de l'equipe ayant acces a la plateforme ouvre dedomena,
tape une question et recoit une reponse. Voici des exemples realistes de ce que
nous pourrions ecrire:

\begin{highlight}
  \textit{"Quels sont nos cinq meilleurs comptes en chiffre d'affaires le trimestre dernier?"}

  \vspace{0.4em}
  \textit{"Donne-moi un resume de tous les tickets d'assistance ouverts ce mois-ci."}

  \vspace{0.4em}
  \textit{"Quelle gamme de produits a le taux de retour le plus eleve selon nos donnees de vente?"}

  \vspace{0.4em}
  \textit{"Compare nos effectifs de janvier a aujourd'hui et signale les departements
           qui ont cru de plus de 20 pour cent."}
\end{highlight}

dedomena lit toutes les sources connectees simultanement, trouve les informations
pertinentes et redige une reponse claire avec des references indiquant d'ou
viennent les donnees.

\subsection{Etape 3 : agir sur les resultats}

La plateforme ne retourne pas seulement un chiffre. Elle explique ce qu'elle a
trouve, signale les anomalies et presente les informations sous forme de resume,
de tableau ou de liste, selon ce qui est le plus utile. La personne qui pose la
question dispose de tout ce dont elle a besoin pour prendre une decision sans
ouvrir un seul tableur.


% ── 4. Qui en beneficie ───────────────────────────────────────────────────────
\section{Qui en beneficie et comment}

\vspace{0.4em}
\begin{center}
\begin{tikzpicture}[
    role/.style={rectangle, rounded corners=6pt,
                 draw=cengine!40, fill=cengine!7,
                 font=\small\bfseries, text=darkgray,
                 minimum width=2.8cm, minimum height=0.85cm,
                 text centered, inner sep=8pt},
    gain/.style={rectangle, rounded corners=4pt,
                 draw=coutput!40, fill=coutput!7,
                 font=\scriptsize, text=darkgray,
                 minimum width=5.4cm, text width=5.2cm,
                 inner sep=7pt},
    conn/.style={-{Stealth[length=5pt]}, thick, color=mutedtext!60}]

  \node[role] (mgr)   at (0,  3.0) {Direction};
  \node[role] (sales) at (0,  1.5) {Ventes};
  \node[role] (ops)   at (0,  0.0) {Operations};
  \node[role] (fin)   at (0, -1.5) {Finances};
  \node[role] (data)  at (0, -3.0) {Analystes};

  \node[gain] (gmgr)  at (7.6,  3.0)
    {Reponses instantanees aux questions de performance. Plus d'attente pour les rapports.};
  \node[gain] (gsales)at (7.6,  1.5)
    {Historique client et sante des comptes disponibles en un instant.};
  \node[gain] (gops)  at (7.6,  0.0)
    {Une vue combinee des stocks, fournisseurs et donnees de processus en un seul endroit.};
  \node[gain] (gfin)  at (7.6, -1.5)
    {Croisement des donnees financieres et detection des depassements sans formules.};
  \node[gain] (gdata) at (7.6, -3.0)
    {Liberes des demandes de donnees routinieres. Concentres sur l'analyse a forte valeur ajoutee.};

  \foreach \r/\g in {mgr/gmgr, sales/gsales, ops/gops, fin/gfin, data/gdata}{
    \draw[conn] (\r) -- (\g);
  }
\end{tikzpicture}
\end{center}
\vspace{0.4em}


% ── 5. Ce qui le distingue ────────────────────────────────────────────────────
\section{Ce qui le distingue des autres outils}

\vspace{0.6em}
\begin{center}
\begin{tikzpicture}[
    col/.style={rectangle, rounded corners=6pt,
                draw=#1!50, fill=#1!8,
                minimum width=5.6cm, text width=5.4cm,
                inner sep=12pt, font=\small, text=#1!80!black}]

  \node[font=\small\bfseries\color{cdanger!80}] at (0, 0)   {Sans dedomena};
  \node[font=\small\bfseries\color{cokay!80}]   at (7.4, 0) {Avec dedomena};

  \draw[thick, dashed, color=mutedtext!40] (3.7, 0.5) -- (3.7, -7.8);

  \foreach \y/\bad/\good in {
    -1.0/{Attendre 1 a 3 jours pour une reponse}/{Reponse en moins de 30 secondes},
    -2.6/{Un tableau de bord a la fois}/{Lit toutes les sources simultanement},
    -4.2/{Necessite SQL ou Excel}/{Questions en francais courant uniquement},
    -5.8/{L'analyste est le goulot d'etranglement}/{N'importe quel membre peut poser une question},
    -7.4/{Donnees dispersees dans 10 outils}/{Un seul endroit pour toutes les questions}
  }{
    \node[col=cdanger, align=left] at (0, \y)   {\bad};
    \node[col=cokay,   align=left] at (7.4, \y) {\good};
    \node[font=\large\color{accent}] at (3.7, \y) {$\rightarrow$};
  }

\end{tikzpicture}
\end{center}
\vspace{0.4em}

\subsection{Les donnees restent chez nous}

C'est essentiel du point de vue de la conformite. dedomena ne telecharge pas vos
donnees brutes vers un serveur externe. Lorsqu'une question est posee, seules les
parties des donnees pertinentes pour cette question specifique sont envoyees a
l'intelligence artificielle pour generer la reponse. Le reste reste en place,
dans notre environnement.


% ── 6. Ce que nous y gagnons ──────────────────────────────────────────────────
\section{Ce que nous y gagnons en tant qu'organisation}

\begin{enumerate}[leftmargin=1.6em, itemsep=8pt]
  \item \textbf{Des decisions plus rapides.}
  Les questions qui prenaient des jours ne prennent plus que quelques secondes.
  Les equipes agissent sur l'information pendant qu'elle est encore d'actualite.

  \item \textbf{Moins de dependance envers le personnel technique pour les questions courantes.}
  Les analystes sont liberes des demandes de donnees basiques.
  Ils concentrent leur temps sur des travaux qui necessitent reellement leurs competences.

  \item \textbf{Un seul endroit pour interroger toutes nos donnees.}
  Au lieu de se connecter a cinq outils differents et d'exporter cinq tableurs,
  tout le monde passe par une seule interface.

  \item \textbf{Un historique complet de ce qui a ete demande et quand.}
  Chaque question est enregistree avec les sources utilisees, creant une piste
  d'audit naturelle pour toute decision basee sur les donnees.

  \item \textbf{Une meilleure culture de la donnee dans l'organisation.}
  Quand les gens peuvent poser des questions directement, ils s'engagent plus
  souvent avec les donnees. Avec le temps, cela construit une culture ou les
  decisions sont fondees sur des preuves.
\end{enumerate}


% ── 7. Plan de deploiement ────────────────────────────────────────────────────
\section{Plan de deploiement suggere}

\vspace{0.8em}
\begin{center}
\begin{tikzpicture}[
    phase/.style={rectangle, rounded corners=6pt,
                  draw=#1!60, fill=#1!10,
                  minimum width=4.0cm, text width=3.8cm,
                  inner sep=12pt, font=\small, text=darkgray,
                  drop shadow={shadow xshift=1pt, shadow yshift=-1pt,
                               fill=#1!20, opacity=0.45}},
    badge/.style={circle, draw=#1!60, fill=#1!15,
                  font=\small\bfseries, text=#1!80!black,
                  minimum size=0.9cm},
    phaseconn/.style={-{Stealth[length=6pt]}, thick, color=mutedtext!50}]

  \node[phase=camber] (p1) at (0, 0) {
    \textbf{Phase 1 : Pilote}\\[4pt]
    Connecter 2 ou 3 sources cles.\\
    Test de 2 semaines avec les directeurs et responsables d'equipe.\\
    Recueillir les retours.
  };

  \node[phase=cengine] (p2) at (5.2, 0) {
    \textbf{Phase 2 : Deploiement}\\[4pt]
    Connecter toutes les sources.\\
    Ouvrir l'acces a toutes les equipes.\\
    Session de formation courte.
  };

  \node[phase=coutput] (p3) at (10.4, 0) {
    \textbf{Phase 3 : Optimisation}\\[4pt]
    Analyser les questions les plus frequentes.\\
    Construire des resumes permanents pour les besoins recurrents.
  };

  \draw[thick, color=mutedtext!40] (-2, -1.8) -- (12.4, -1.8);

  \node[badge=camber]  at (0,    -1.8) {1};
  \node[badge=cengine] at (5.2,  -1.8) {2};
  \node[badge=coutput] at (10.4, -1.8) {3};

  \node[font=\scriptsize\color{mutedtext}] at (0,    -2.55) {Semaines 1 et 2};
  \node[font=\scriptsize\color{mutedtext}] at (5.2,  -2.55) {Mois 2};
  \node[font=\scriptsize\color{mutedtext}] at (10.4, -2.55) {Mois 3 et au-dela};

  \draw[phaseconn] (p1) -- (p2);
  \draw[phaseconn] (p2) -- (p3);

\end{tikzpicture}
\end{center}
\vspace{0.4em}


% ── 8. En resume ──────────────────────────────────────────────────────────────
\section{En resume}

dedomena donne a notre organisation la capacite de poser n'importe quelle question
sur nos donnees et d'obtenir une reponse claire et fiable en quelques secondes,
sans ecrire de code, sans attendre un rapport, et sans avoir besoin de savoir ou
se trouvent les donnees.

Nous allons l'utiliser pour reduire le temps entre une question et une decision,
pour supprimer le goulot d'etranglement qui repose aujourd'hui sur notre equipe
d'analyse, et pour donner a chaque membre de l'equipe un acces direct aux
informations dont il a besoin pour bien faire son travail.

\begin{callout}
  \textbf{L'objectif est simple :} n'importe qui dans l'organisation doit pouvoir
  ouvrir dedomena, taper une question sur notre activite, et avoir une reponse
  claire devant lui avant que son cafe ne refroidisse.
\end{callout}

\vspace{2em}
\textcolor{mutedtext}{\small Ce document est destine a un usage interne. Pour une
demonstration en direct ou pour discuter de l'integration avec nos outils existants,
veuillez contacter l'equipe produit.}

\end{document}
